% =====================================================================
%  config/dados-da-equipe.tex
%  ESTE É O ÚNICO ARQUIVO DE CONFIGURAÇÃO QUE A EQUIPE PRECISA EDITAR.
%  Leva dois minutos. Faça isso antes de escrever qualquer coisa.
% =====================================================================

% --- 1. Qual ciclo de projeto este relatório documenta? ---------------
%     1 = AP1 (sensoriamento e atuação)
%     2 = AP2 (conectividade e conjunto de dados)
%     3 = AP3 (projeto integrador de AIoT)
%     Trocar este número liga e desliga capítulos inteiros do relatório.
\def\CICLO{2}

% --- 2. Mostrar as instruções cinzas ("Como preencher")? --------------
%     Mantenha ligado enquanto escreve.
%     COMENTE a linha abaixo para gerar a versão de entrega.
%     Versão de entrega (25/10): comentada, as caixas cinzas somem.
% \def\MOSTRARNOTAS{}

% --- 3. Identificação do projeto --------------------------------------
\newcommand{\tituloProjeto}{TermoLOG}
\newcommand{\subtituloProjeto}{Monitoramento e sinalização local de temperatura e umidade em ambientes críticos}
\newcommand{\nomeEquipe}{CMDBD Soluções}
\newcommand{\repositorioEquipe}{https://github.com/victor-kauan-coder/TermoStats}

% --- 4. Integrantes ---------------------------------------------------
%     Uma linha por integrante: \integrante{Nome completo}{Matrícula}{Papel no ciclo}
%     Papéis do Módulo 1: integração, firmware, hardware, documentação.
%     Papéis do Módulo 2: integração, firmware e comunicação, dados
%     (gravador, campanha e qualidade), documentação. Os papéis rotacionam:
%     ninguém repete na AP2 o papel que teve na AP1.
%
%     Papéis na AP1 (para conferir a rotação):
%       João Pedro = integração; Victor e Davi = firmware;
%       Felipe = hardware/testes; Luanderson = documentação.
%     Rotação confirmada pela equipe em 25/10. Os commits da AP2 foram
%     enviados por um único integrante; a autoria das tarefas está no Cap. 9.
\newcommand{\listaIntegrantes}{%
  \integrante{Joao Pedro Saleh}{20249012194}{dados (gravador e qualidade)}
  \integrante{Victor Kauan da Silva Miranda}{20249021999}{integração}
  \integrante{Davi Sousa Soares}{20249022000}{documentação}
  \integrante{Felipe Gabriel Soares Rodrigues}{20249020884}{firmware e comunicação}
  \integrante{Luanderson Oliveira de Carvalho}{20249012050}{dados (campanha e registro de eventos)}
}

% --- 5. Dados institucionais (normalmente não mudam) ------------------
\newcommand{\instituicaoCurso}{Universidade Federal do Piauí}
\newcommand{\unidadeCurso}{Centro de Ciências da Natureza}
\newcommand{\programaCurso}{Departamento de Computação}
\newcommand{\nomeDisciplina}{Tópicos em Sistemas Computacionais I --- Internet das Coisas e AIoT}
\newcommand{\periodoLetivo}{2026.2}
\newcommand{\docenteDisciplina}{Prof. Dr. Ricardo de Andrade Lira Rabelo \textbullet\ Ministrante: Atos Apollo Silva Borges}
\newcommand{\localEntrega}{Teresina}
\newcommand{\dataEntrega}{25 de outubro de 2026}

% --- 6. Uso de IA generativa ------------------------------------------
%     A disciplina exige declaração, não proíbe o uso.
%     Preencha o Apêndice D; aqui vai só o resumo de uma linha.
\newcommand{\resumoUsoIA}{%
  A equipe usou IA generativa (Claude, da Anthropic) apenas como apoio:
  consulta de erros, auxílio na escrita do código e auxílio na montagem
  do relatório. Ideias, decisões de projeto, testes e coleta dos dados são
  da equipe. O detalhamento fica no Apêndice~\ref{ap:ia}.}
